\chapter*{How to read this book}
\addcontentsline{toc}{chapter}{How to read this book}

Read this first: IAM is not a new model of gravity or of reality. It is general relativity with one piece added to the entropy Jacobson
already wrote down, and that piece is information.

\section*{Status labels}
Every result carries a label. A reader who remembers nothing else should remember which numbers are derived and which were set from data.

\begin{center}\small
\begin{tabularx}{\textwidth}{@{}lX@{}}
\toprule
\derived & Follows from stated premises by algebra or standard physics, with no parameter adjusted to data. \\
\calc & A number computed from a derived formula and published constants or measurements; recomputed for this edition. \\
\calibrated & A value set once from healthy or reference data for one system and then held fixed. The data are named. \\
\measured & A value read from data by the instruments of this book and frozen in a file in the repository. \\
\observed & A published measurement, quoted with its source. \\
\fitted & A parameter adjusted to make a formula match the quantity it is later compared with. The agreement is then not evidence. \\
\conjecture & A structural or interpretive claim that is not yet a calculation. \\
\analogy & A correspondence between two domains that guides a method but is not a derivation. \\
\interp & How we read a set of results; it adds no result of its own. \\
\prediction & A falsifiable statement about a measurement not yet made, with the test named. \\
\openprob & A known gap or inconsistency, stated openly. \\
\bottomrule
\end{tabularx}
\end{center}

\section*{The path through the book}
Part~\ref{part:1} is required reading for everything after it. Each of the next five Parts belongs to one field and can be reviewed on its own
by a specialist in that field, with Part~\ref{part:1}:
\begin{itemize}\itemsep0pt
\item a cosmologist: Part~\ref{part:2};
\item a relativist or quantum-gravity theorist: Part~\ref{part:bh};
\item a particle physicist or a specialist in the foundations of quantum mechanics: Part~\ref{part:qp};
\item a quantum-hardware or device physicist: Part~\ref{part:3};
\item a geneticist or biological physicist: Part~\ref{part:4}, with the one-gauge chapter (Chapter~\ref{ch:onegauge}).
\end{itemize}
Part~\ref{part:5} joins them, collects every prediction with its test, and lists what is open; it assumes all of them.

\section*{What this book does not contain}
Instrument software beyond what is needed to reproduce a number, calibration tooling and commercial material are not reproduced. The book gives no
medical, nutritional or treatment advice: the cell instrument measures a state; it does not diagnose, and it is not used on any question until it
has been commissioned for it.

Pick the Part that touches your field, read the derivation it rests on, then open the repository and run it yourself. Every result tells you how strongly it is claimed, so you know exactly where to push.

\section*{Conventions}
SI units throughout, CODATA~2018 constants~\cite{CODATA2018} (Appendix~\ref{app:canon}). Cosmological parameters are Planck~2018
TT,TE,EE+lowE+lensing~\cite{Planck2018VI} unless stated. $\ln$ is the natural logarithm; one bit is $\ln2$ nats. The Mahaffey number is
$\mathcal M=E_{\rm drive}/\kB T$; for one ATP at 37\,\textdegree C it is 20.94, which is 30.2 in Landauer units $\kB T\ln2$. Every number and figure
can be rebuilt from the repository (Appendix~\ref{app:repro}).
